\section{Method}
\label{sec:method}

\subsection{Preliminaries}
\label{sec:prelim}

\textbf{Low-rank adaptation (LoRA).}
LoRA~\cite{hu2022lora} freezes a pre-trained weight matrix $W_0\in\mathbb{R}^{d_{\text{out}}\times d_{\text{in}}}$
and injects a trainable low-rank residual,
\begin{equation}
W \;=\; W_0 \;+\; \Delta W \;=\; W_0 + \frac{\alpha}{r}\,B A,
\label{eq:lora}
\end{equation}
where $A\in\mathbb{R}^{r\times d_{\text{in}}}$, $B\in\mathbb{R}^{d_{\text{out}}\times r}$
and $r\ll \min(d_{\text{in}},d_{\text{out}})$. The scalar $\alpha/r$ absorbs the rank
normalization. The forward pass of an adapted layer is
\begin{equation}
h \;=\; W_0 x + \frac{\alpha}{r}\,B A x \;=\; W_0 x + \delta,
\qquad \delta := \frac{\alpha}{r}\,B A x \in \mathbb{R}^{d_{\text{out}}},
\label{eq:forward}
\end{equation}
where $\delta$ denotes the adapter's residual contribution.

\textbf{Class-incremental continual learning (CIL).}
We consider a sequence of $T$ tasks $\{\mathcal{T}_t\}_{t=1}^{T}$. Task $\mathcal{T}_t$
provides a training set $\mathcal{D}_t=\{(x_i,y_i)\}$ over $c$ new classes, and the
classifier head is grown to cover all classes seen so far. At test time the task
identity is not available, so prediction is made over the union of all seen classes.
A single shared LoRA adapter is trained sequentially; its parameter count is
$r(d_{\text{in}}+d_{\text{out}})$ per layer and is \emph{independent} of $T$.
% NOTE (2026-09-24): the class count is written $c$, not $k$. $k$ is reserved
% throughout this paper for the number of protected directions (the rank budget of
% the importance kernel); writing "over $k$ new classes" here put it within a few
% lines of "top-$k$ eigendirections" and "budget $k$" with a different meaning.
% Do not reintroduce $k$ for the class count.

\subsection{Motivation: forgetting as subspace overlap}
\label{sec:motivation}

Let $\xi_\tau$ collect all \emph{trainable} parameters at the end of task
$\mathcal{T}_\tau$ --- the adapter parameters $\theta_\tau$ (the flattened low-rank factors
of Sec.~\ref{sec:folora}) together with the classifier-head parameters --- and let
$\delta\xi=\xi-\xi_\tau$ be a later task's displacement of them. A second-order Taylor
expansion of task $\tau$'s loss around $\xi_\tau$ gives
\begin{equation}
\mathcal{L}_\tau(\xi_\tau+\delta\xi)
\;\approx\; \mathcal{L}_\tau(\xi_\tau)
+ \nabla_\xi\mathcal{L}_\tau^{\!\top}\delta\xi
+ \tfrac12\, \delta\xi^{\!\top} \mathcal{F}_\tau\, \delta\xi,
\label{eq:taylor}
\end{equation}
where $\mathcal{F}_\tau=\mathbb{E}_{(x,y)\sim\mathcal{D}_\tau}\big[\nabla_\xi \log p(y|x)\,\nabla_\xi\log p(y|x)^{\!\top}\big]$
is the empirical Fisher information matrix of task $\tau$ over those parameters.
% NOTE (2026-09-24): Eq. (3) previously read L_tau(theta_tau + delta W) with the linear term
% <grad_{delta W}, delta W> and the quadratic delta W^T F_tau delta W -- adding a weight
% MATRIX to a parameter VECTOR and contracting that matrix with a Fisher defined over the
% vector. It is now in xi (the full trainable vector), exactly as Sec. 4, and it carries the
% CALLIGRAPHIC Fisher because that is the full-parameter object; roman $F_\tau$ (Eq. 5) is
% its adapter block, and is what FOLoRA actually estimates. Do not let the two drift apart:
% a reviewer reads Secs. 3 and 4 back to back, and a symbol that changes meaning between
% them reads as a mistake even when each use is individually defensible.
At the end of task $\tau$ the parameters are (approximately) stationary for that task's
loss, $\nabla_\xi\mathcal{L}_\tau(\xi_\tau)\approx 0$; this is stated as an explicit
assumption (Assumption~\ref{ass:stationary}) and its residual under the regularizer is
quantified in Sec.~\ref{sec:theory}. Hence the \emph{forgetting} of task $\tau$
--- the increase in its loss caused by subsequent learning --- is dominated by the
quadratic term $\frac12\,\delta\xi^{\!\top}\mathcal{F}_\tau\,\delta\xi$. Forgetting is
therefore driven by the alignment of the new displacement with the directions in
which task $\tau$'s loss is sharply curved (i.e., the top eigendirections of
$\mathcal{F}_\tau$). Of that quadratic form, FOLoRA acts on the part involving the adapter
parameters: its regularizer penalizes alignment of the adapter displacement $\delta\theta$
with the directions in which the adapter block of $\mathcal{F}_\tau$ is sharply curved.
Sec.~\ref{sec:theory} makes this precise, including the head terms that FOLoRA does not
constrain.

This observation separates FOLoRA from prior work: O-LoRA~\cite{wang2023olora}
treats every direction of the past low-rank subspace as equally important, whereas
Eq.~\eqref{eq:taylor} shows that directions should be protected \emph{in proportion to
their curvature}. Protecting all directions equally over-constrains the adapter and
wastes the fixed rank budget on directions that carry little information; weighting by
curvature spends that budget where forgetting actually happens. We note that this is an
argument about \emph{where the budget is spent}, not a claim that weighting wins: the
theory does not predict an accuracy advantage for the weighted variant
(Sec.~\ref{sec:equalweight} states the caveat explicitly), and Sec.~\ref{sec:ablation}
reports what our experiments show.

\subsection{FOLoRA}
\label{sec:folora}

We now describe FOLoRA. For each adapted layer, let $\theta\in\mathbb{R}^{d}$ denote the
flattened adapter parameters (the LoRA matrices $A$ and $B$, so $d=r(d_{\text{in}}+d_{\text{out}})$),
and let $g_{\theta}=\nabla_{\theta}\log p(y|x)$ denote the gradient of the log-likelihood
with respect to these parameters.

\paragraph{Importance estimation}
After finishing task $\mathcal{T}_\tau$, we estimate a per-layer \emph{importance
kernel} from the per-sample parameter-gradient autocorrelation,
\begin{equation}
F_\tau \;=\; \mathbb{E}_{(x,y)\sim\mathcal{D}_\tau}\big[\,g_{\theta}(x,y)\,g_{\theta}(x,y)^{\!\top}\,\big]
\;\in\;\mathbb{R}^{d\times d},
\label{eq:kernel}
\end{equation}
% NOTE (2026-09-24): the expectation is over the joint (x,y), as in Eq. (3). It previously
% appeared as E_{x~D_tau}[g(x)g(x)^T], which is not a well-defined gradient -- g depends on
% the label through log p(y|x). This and Eq. (3) were the only two sites of that slip; if
% you add another E[...g(x)...], write the label in.
which is the empirical Fisher information of task $\tau$ restricted to the adapter
parameters --- i.e.\ the adapter block $\mathcal{F}_\tau^{\theta\theta}$ of the
full-parameter Fisher of Sec.~\ref{sec:theory} --- and which we write $F_\tau$ for short
below. To keep the estimate tractable we never materialize the full $d\times d$
matrix; we collect the per-sample gradients into a matrix (one row per sample) and retain
only its top-$k$ singular directions, i.e., a rank-$k$ approximation of $F_\tau$. We
accumulate these kernels across tasks,
\begin{equation}
\bar{F} \;=\; \sum_{\tau < t} F_\tau.
\label{eq:acc}
\end{equation}
Eq.~\eqref{eq:acc} is the \emph{object} the protected subspace is meant to represent.
The implementation stores, per layer, a rank-$k$ factor rather than the full $F_\tau$, so
what is actually accumulated is computed on the concatenation of the retained factors of
all previous tasks; this is an incremental-PCA approximation to the sum in
Eq.~\eqref{eq:acc} rather than the sum itself, and it is exact only when the retained
top-$k$ subspace of the concatenated factors coincides with the union of the per-task
top-$k$ subspaces. We keep the exact form in the equations for readability and use the
approximation in all experiments; the two agree in the limit $k\to d$.
% NOTE (2026-09-24): do not delete this paragraph to make the method look cleaner. The code
% (methods/folora_v2.py) does acc_grads[i] = _topk_directions(cat([acc_grads[i], G]), k),
% i.e. it re-truncates to k after each task, so a direction that ranks just outside the top
% k for task 1 can never be recovered by task 2 -- the accumulation is genuinely lossy.
% Claiming "F_bar = sum F_tau" with no qualifier is the kind of small unfaithfulness a
% reviewer who reads the code will find, and it costs more than it buys.

\paragraph{Protected subspace}
At the start of task $\mathcal{T}_t$, we eigendecompose the accumulated kernel
$\bar{F}=\sum_j \sigma_j u_j u_j^{\!\top}$, where $\{\sigma_j\}_{j}$ are the
eigenvalues sorted in decreasing order and $\{u_j\}$ the corresponding orthonormal
eigenvectors in $\mathbb{R}^{d}$. We retain the top-$k$ eigendirections as
an \emph{importance-weighted protected subspace}
$\mathcal{P}_t=\{(\sigma_j,u_j)\}_{j=1}^{k}$.

\paragraph{Orthogonalization regularizer}
Let $\theta^{t-1}$ denote the adapter parameter snapshot recorded after task
$\mathcal{T}_{t-1}$, and let $\delta\theta=\theta-\theta^{t-1}$ be the parameter
displacement accumulated while training task $\mathcal{T}_t$. During the training of task
$\mathcal{T}_t$, FOLoRA minimizes
\begin{equation}
\mathcal{L} \;=\; \mathcal{L}_{\text{CE}}(\theta) \;+\; \lambda\,
\underbrace{\sum_{j=1}^{k}\sigma_j\,\big(u_j^{\!\top}\,\delta\theta\big)^{2}}_{\Omega(\delta\theta)},
\label{eq:loss}
\end{equation}
where $\lambda>0$ controls the strength and each $u_j\in\mathbb{R}^{d}$ is a protected
direction in the same parameter space as $\delta\theta$. The term
$(u_j^{\!\top}\delta\theta)^2$ penalizes the alignment of the \emph{displacement}
$\delta\theta=\theta-\theta^{t-1}$---rather than the absolute $\theta$---with the
protected direction $u_j$, weighted by the importance $\sigma_j$. Using the displacement
is essential: at the start of task $\mathcal{T}_t$ we have $\delta\theta=0$, so the
regularizer vanishes and the adapter's existing alignment with past tasks is preserved;
only \emph{new} movement into protected directions is penalized. Penalizing the absolute
$\theta$ would instead push the adapter away from the past tasks' directions and
actively erase them, defeating the purpose of protection. Since the regularizer is a
differentiable function of $A$ and $B$, it integrates with standard backpropagation and
requires no hard projection, no replay buffer, and no task identity at inference.

\paragraph{Why importance weighting matters}
Setting all $\sigma_j=1$ reduces $\Omega$ to an equal-weight orthogonalization, which is
the geometry underlying O-LoRA-style orthogonal subspace methods
(Sec.~\ref{sec:equalweight} states the relation exactly). In contrast, FOLoRA's
eigenvalue weighting (i) strongly suppresses reuse of high-curvature directions, which is
where forgetting is concentrated, and (ii) leaves low-curvature directions nearly free, so
the adapter can reuse them for new tasks without wasting rank. This is the mechanism by
which FOLoRA spends a fixed rank budget on the directions the bound identifies as costly.
We note that the theory does not by itself predict an accuracy advantage for weighting
over equal weighting --- a larger penalty family is not automatically a better one, and
Sec.~\ref{sec:equalweight} says so explicitly; whether it helps in practice is an
empirical question, and Sec.~\ref{sec:ablation} reports what our experiments show.

% ============================================================================
% Fig. 1 -- FOLoRA framework.  \input from 03_method.tex.
%
% Requires (already loaded in main.tex):
%   \usepackage{tikz}
%   \usetikzlibrary{arrows.meta,positioning,fit,backgrounds,calc}
%
% Layout: two rows. The solid row is the forward pass of one task, left to right.
% The dashed row is the one-time update at the end of a task, drawn right to left
% because it feeds back into the top row -- the picture is a loop, not a pipeline.
%
% GEOMETRY IS LOAD-BEARING. Compiled with elsarticle [review]: the text block is
% only 390pt x 548.5pt and \topfraction is 0.7, i.e. a top float may occupy at most
% 384pt (~135mm) INCLUDING its caption. The picture below plus its caption is
% ~115mm. If you lengthen the caption or move a node down, recompile and check both
% (a) "Overfull \hbox ... too wide" and (b) that the figure still lands mid-paper --
% an over-tall float silently defers to a float page after the bibliography.
%
% The explicit \linespread{1}\selectfont in the node style is also load-bearing:
% [review] is double-spaced, and without it every baselineskip inside a node is
% scaled by ~1.67 and the boxes grow into each other.
%
% There is deliberately no numeric content here. Every table and every other
% figure in the paper carries measured numbers; this one carries only the
% architecture of the method, so it does not need re-deriving when the numbers are.
% ============================================================================
\begin{figure}[t]
\centering
\begin{tikzpicture}[
  x=1mm, y=1mm,
  every node/.append style={font=\small,
    execute at begin node={\linespread{1}\selectfont}},
  box/.style={draw, rounded corners=1.5pt, align=center, inner sep=3pt,
              minimum height=11mm, text width=24mm},
  frozen/.style={box, fill=black!7},
  trainable/.style={box, fill=black!18},
  state/.style={box, fill=white, dashed},
  ar/.style={-{Latex[length=1.7mm]}, semithick},
  dar/.style={-{Latex[length=1.7mm]}, semithick, dashed},
]

% ---------- solid row: the forward pass of one task ----------
\node (in)   at (  -6, 0) {$x$};
\node (bb)   at (  14, 0) [frozen, text width=20mm]
  {frozen\\[-1pt]\scriptsize ViT-B/16};
\node (ad)   at (  46, 0) [trainable, text width=22mm]
  {LoRA $\theta$\\[-1pt]\scriptsize rank $r$, shared};
\node (hd)   at (  74, 0) [trainable, text width=22mm]
  {head $\psi$\\[-1pt]\scriptsize dropped at test};
\node (loss) at ( 106, 0) {$\mathcal{L}_{\mathrm{CE}}$};

\draw[ar] (in)   -- (bb);
\draw[ar] (bb)   -- (ad);
\draw[ar] (ad)   -- (hd);
\draw[ar] (hd)   -- (loss);

% ---------- dashed row: once, at the end of a task, right to left ----------
\node (fis) at ( 106, -30) [state, text width=32mm]
  {empirical Fisher\\[-1pt]\scriptsize of $\theta$, per layer};
\node (eig) at (  70, -30) [state, text width=30mm]
  {top-$k$\\ eigendirections\\[-1pt]\scriptsize $\{(\sigma_j,u_j)\}_{j\le k}$};
\node (om)  at (  29, -30) [state, text width=38mm]
  {$\Omega(\delta\theta)$\\[-1pt]\scriptsize
   $\sum_{j\le k}\sigma_j\,(u_j^{\!\top}\delta\theta)^{2}$};

\draw[dar] (loss) -- node[left=2pt, align=right]
  {\scriptsize once, at the\\[-1pt]\scriptsize end of task $t$} (fis);
\draw[dar] (fis)  -- (eig);
\draw[dar] (eig)  -- (om);

% ---------- feedback: the penalty acts on the adapter ----------
% Elbow path, not the -| syntax: -| would put its last segment along ad's own
% bottom edge, which reads as a line stuck to the box.
\draw[dar] (om.north) -- (29, -13) -- (46, -13) -- (ad.south);
\node at (48, -15.5) [anchor=west] {\scriptsize $\lambda\,\Omega$};

% ---------- the property the figure exists to show ----------
\node at (57, -46) [align=center, text width=110mm]
  {\scriptsize One adapter for all tasks: its trainable parameter count does not
   grow with the number of tasks. Only the importance state
   $\{(\sigma_j,u_j)\}$ and the reference $\theta^{t-1}$ are updated between tasks.};

\end{tikzpicture}
\caption{FOLoRA. \textbf{Solid row}: the forward pass of task $t$, where a single shared
rank-$r$ adapter $\theta$ is trained on a frozen ViT-B/16 backbone with the penalty
$\lambda\Omega$ added to the cross-entropy loss; the head $\psi$ is trained but discarded,
since prediction is by nearest class mean in the frozen feature space.
\textbf{Dashed row}: once, after task $t$ ends, the adapter's empirical Fisher is
re-estimated per layer, its top $k$ eigendirections are extracted, and they are
accumulated into the kernel $\bar{F}$ defining $\Omega$ for the next task. The adapter's
rank, hence its trainable parameter count, is fixed for the whole sequence; only the
importance state grows, so the memory cost is an importance-model cost, not a parameter
cost (Table~\ref{tab:resources}).}
\label{fig:framework}
\end{figure}


Fig.~\ref{fig:framework} shows the procedure end to end: the forward pass of one task on
the solid row, and the one-time update that produces the penalty for the next task on the
dashed row. Two properties of that picture are the ones the rest of the paper relies on:
the adapter and its parameter count are shared across all tasks, so nothing in the solid
row grows with $T$; and the quantity that is updated between tasks is the importance state
$\{(\sigma_j,u_j)\}$, so the memory FOLoRA spends is an importance-model cost rather than a
parameter cost (quantified in Table~\ref{tab:resources}). The same procedure is given in
pseudocode as Algorithm~\ref{alg:folora}.

\begin{algorithm}[t]
\caption{FOLoRA for class-incremental continual learning}
\label{alg:folora}
\begin{algorithmic}[1]
\STATE \textbf{Input:} frozen backbone; LoRA rank $r$; strength $\lambda$; budget $k$.
\STATE Initialize shared LoRA adapters $\Delta W^{(l)}=\frac{\alpha}{r}B^{(l)}A^{(l)}$ for each adapted layer $l$.
\STATE $\bar{F}^{(l)} \leftarrow \mathbf{0}$ for all $l$.
\FOR{each task $\mathcal{T}_t$, $t=1,\dots,T$}
  \STATE \textbf{if} $t>1$: eigendecompose $\bar{F}^{(l)}=\sum_j\sigma_j u_j u_j^{\!\top}$,
         keep top-$k$ $\{(\sigma_j,u_j)\}$.
  \STATE Train adapter parameters $\theta^{(l)}$ and the grown classifier head by minimizing
         $\mathcal{L}_{\text{CE}}+\lambda\sum_{l}\sum_{j}\sigma_j\big(u_j^{\!\top}\big(\theta^{(l)}-\theta^{(l)}_{t-1}\big)\big)^2$.
  \STATE Record adapter snapshot $\theta^{(l)}_{t}\leftarrow\theta^{(l)}$.
  \STATE Estimate $F_t^{(l)}=\mathbb{E}[g_{\theta}\,g_{\theta}^{\!\top}]$ on $\mathcal{D}_t$; update
         $\bar{F}^{(l)}\leftarrow\bar{F}^{(l)}+F_t^{(l)}$.
\ENDFOR
\STATE \textbf{Output:} a single shared adapter with fixed parameter budget.
\end{algorithmic}
\end{algorithm}

\subsection{Relation to prior methods}
\label{sec:relation}
\begin{itemize}
\item \textbf{EWC}~\cite{kirkpatrick2017overcoming} penalizes per-parameter quadratic
displacement $\sum_i F_{ii}(\theta_i-\theta_i^*)^2$ using the \emph{diagonal} Fisher.
FOLoRA instead uses the adapter block of the second-order expansion and applies it as a
rank-$k$ subspace projection rather than a pointwise penalty. The distinction is not that
one is exact and the other is not --- both truncate the curvature, EWC to its diagonal and
FOLoRA to its top-$k$ spectrum --- but that they discard different things: EWC keeps all
$d$ coordinates and drops every correlation among them, whereas FOLoRA keeps the $k$
directions of largest curvature with their correlations intact and drops the rest. Which
truncation suits a fixed-rank adapter is an empirical question, answered for this setting
by the comparison in Sec.~\ref{sec:main_results}.
\item \textbf{O-LoRA}~\cite{wang2023olora} allocates a new adapter per task and
initializes it orthogonally to previous tasks' input subspaces, treating all directions
equally; FOLoRA keeps a \emph{single} fixed-budget adapter and protects directions in
proportion to their Fisher importance.
\item \textbf{GPM}~\cite{saha2021gradient} and related null-space methods project
gradients onto a null space; FOLoRA realizes a \emph{soft}, importance-weighted
orthogonal projection that avoids the rank saturation of hard projections.
\end{itemize}
